\section{Sub-Shot-Noise Imaging}
\label{sec:sub_shot_noise}

Sub-shot-noise imaging exploits quantum correlations between photon pairs to achieve sensitivity below the classical shot-noise limit, which bounds the minimum noise in measurements made with coherent (classical) light sources.

\subsection{Quantum Noise Limits in Imaging}

\citet{brida2010} provided the first experimental demonstration of sub-shot-noise quantum imaging, achieving a noise reduction of approximately 40\% below the shot-noise level using spatially entangled photon pairs from SPDC. \citet{samantaray2017} subsequently realised the first sub-shot-noise wide-field microscope, extending the technique from proof-of-concept to a practical imaging modality capable of operating over extended fields of view.

The quantum enhancement in these systems arises from the photon-number correlations between signal and idler beams: by subtracting the idler intensity from the signal, classical noise sources are suppressed below the shot-noise floor. However, the improvement is limited by optical losses and detector inefficiencies, which degrade the quantum correlations.

\subsection{Deep Learning Denoising for Quantum-Enhanced SNR}

% add: noise models (Poissonian vs sub-Poissonian), DL denoising archs, compare w/ BM3D/NLM

Deep learning denoising represents a promising direction for extending the practical advantage of sub-shot-noise imaging. Classical denoising algorithms assume Gaussian or Poissonian noise statistics, which do not accurately describe the sub-Poissonian noise in quantum-correlated imaging. Neural networks trained on data with appropriate quantum noise models could exploit the specific statistical structure of quantum measurements to achieve superior denoising performance.

Furthermore, deep learning may help compensate for the optical losses that currently limit practical sub-shot-noise imaging, effectively widening the parameter regime in which quantum advantage is achievable. This possibility has been explored in related contexts such as deep learning microscopy \citep{rivenson2017}, but dedicated studies in the quantum imaging setting remain sparse.
